% =====================================================================
%  Phase 1-4 Completion Summary
%  LaTeX rendering of docs/PHASE1-4_COMPLETION_SUMMARY.md
%
%  Compile with:  pdflatex PHASE1-4_COMPLETION_SUMMARY.tex   (run twice,
%                 so the table of contents resolves)
% =====================================================================
\documentclass[11pt,a4paper]{article}

% ---------- encoding & fonts ----------
\usepackage[T1]{fontenc}
\usepackage[utf8]{inputenc}
\usepackage{lmodern}
\usepackage{textcomp}
\usepackage{amssymb}
\usepackage{microtype}

% ---------- page ----------
\usepackage[a4paper,left=2.3cm,right=2.3cm,top=2.5cm,bottom=2.5cm]{geometry}
\usepackage{fancyhdr}
\usepackage{parskip}

% ---------- tables ----------
\usepackage{booktabs}
\usepackage{tabularx}
\usepackage{xltabular}   % longtable + X columns, for tables that span pages
\usepackage{array}

% ---------- colour, symbols, code, links ----------
\usepackage{xcolor}
\usepackage{pifont}
\usepackage{listings}
\usepackage{enumitem}
\usepackage{titlesec}
\usepackage{hyperref}

% ---------------------------------------------------------------------
%  Palette
% ---------------------------------------------------------------------
\definecolor{okgreen}{HTML}{1B7F3B}
\definecolor{nored}{HTML}{B3261E}
\definecolor{warnamber}{HTML}{B36B00}
\definecolor{deferblue}{HTML}{2B5EA8}
\definecolor{rulecol}{HTML}{BFBFBF}
\definecolor{codebg}{HTML}{F5F5F2}
\definecolor{codeframe}{HTML}{D6D6D0}
\definecolor{linkcol}{HTML}{1F4E79}

\hypersetup{
  colorlinks=true,
  linkcolor=linkcol,
  urlcolor=linkcol,
  pdftitle={Phase 1-4 Completion Summary},
  pdfsubject={Unified Ransomware Detection \& Recovery System}
}

% ---------------------------------------------------------------------
%  Status symbols  (stand in for the markdown emoji)
% ---------------------------------------------------------------------
\newcommand{\ok}{\textcolor{okgreen}{\ding{51}}}
\newcommand{\no}{\textcolor{nored}{\ding{55}}}
\newcommand{\warn}{\textcolor{warnamber}{$\blacktriangle$}}
\newcommand{\defer}{\textcolor{deferblue}{\textbf{\textbardbl}}}

% ---------------------------------------------------------------------
%  Column types
% ---------------------------------------------------------------------
\newcolumntype{L}{>{\raggedright\arraybackslash}X}
\newcolumntype{Y}[1]{>{\hsize=#1\hsize\raggedright\arraybackslash}X}
\newcolumntype{P}[1]{>{\raggedright\arraybackslash}p{#1}}
\newcolumntype{K}[1]{>{\centering\arraybackslash}p{#1}}

\renewcommand{\arraystretch}{1.25}
\setlength{\tabcolsep}{5pt}

% ---------------------------------------------------------------------
%  Headings - numbered from 0, matching the markdown's own numbering
% ---------------------------------------------------------------------
\setcounter{secnumdepth}{3}
\setcounter{tocdepth}{2}

\titleformat{\section}
  {\normalfont\Large\bfseries}{\thesection.}{0.6em}{}
  [\vspace{1pt}{\color{rulecol}\titlerule[0.8pt]}]
\titlespacing*{\section}{0pt}{2.4ex plus 1ex minus .2ex}{1.4ex plus .2ex}

\titleformat{\subsection}
  {\normalfont\large\bfseries}{\thesubsection}{0.6em}{}
\titlespacing*{\subsection}{0pt}{2.2ex plus 1ex minus .2ex}{1.0ex plus .2ex}

% ---------------------------------------------------------------------
%  Code blocks
% ---------------------------------------------------------------------
\lstdefinestyle{shell}{
  basicstyle=\ttfamily\small,
  backgroundcolor=\color{codebg},
  frame=single,
  rulecolor=\color{codeframe},
  framesep=7pt,
  framexleftmargin=3pt,
  breaklines=true,
  breakatwhitespace=true,
  columns=fullflexible,
  keepspaces=true,
  showstringspaces=false,
  aboveskip=1.2\baselineskip,
  belowskip=1.2\baselineskip
}
\lstset{style=shell}

% ---------------------------------------------------------------------
%  Running head
% ---------------------------------------------------------------------
\pagestyle{fancy}
\fancyhf{}
\renewcommand{\headrulewidth}{0.4pt}
\renewcommand{\headrule}{\hbox to\headwidth{\color{rulecol}\leaders\hrule height \headrulewidth\hfill}}
\fancyhead[L]{\footnotesize Phase 1--4 Completion Summary}
\fancyhead[R]{\footnotesize\thepage}

\begin{document}

% =====================================================================
%  Front matter
% =====================================================================
\begin{center}
  {\LARGE\bfseries Phase 1--4 Completion Summary}\\[4pt]
  {\large every requirement, checked against the reference PDF}
\end{center}

\vspace{0.8em}

\noindent
\begin{tabularx}{\linewidth}{@{}P{2.4cm}L@{}}
\toprule
\textbf{Reference}  & \emph{Unified Ransomware Detection \& Recovery System --- Complete Project
                      Documentation}, v1.6, 31 January 2026, 68 pages (all 68 extracted and read). \\
\textbf{Repository} & \texttt{D:\textbackslash{}Unified\_Ransomware\_Project}, branch
                      \texttt{fix/phase1-4-remediation} \\
\textbf{Date}       & 12 August 2026 \\
\textbf{Scope}      & Phases 1--4 (Semester 1, Weeks 1--16). Weeks 17--32 items are judged
                      against the document's own timeline, not counted as gaps. \\
\bottomrule
\end{tabularx}

\vspace{1.0em}

\noindent\textbf{Status symbols.}\quad
\ok\ met \quad\textbullet\quad
\warn\ deviation, documented \quad\textbullet\quad
\defer\ deferred to a later phase \quad\textbullet\quad
\no\ not built

\vspace{1.2em}
\hrule height 0.4pt
\vspace{0.6em}

\tableofcontents

\clearpage

% =====================================================================
\setcounter{section}{-1}
\section{Verdict}
% =====================================================================

\noindent
\begin{tabularx}{\linewidth}{@{}P{4.3cm}L@{}}
\toprule
\textbf{Test suite}          & \textbf{371 passed, 2 skipped, 0 failed} (baseline at session start: 305/2/0) \\
\textbf{API contract checks} & \textbf{53 / 53} --- the document's own listings replayed verbatim against the running services \\
\textbf{Table 5.9 benchmarks}& \textbf{10 / 10} met, all re-measured today \\
\textbf{Table 5.8 test cases}& \textbf{11 of 11} in-scope pass, each with a \texttt{tc*}-named test. TC-12 is Weeks 17--24 \\
\textbf{Documented endpoints}& \textbf{10 / 10} implemented, field-for-field \\
\textbf{Blocking defects}    & \textbf{0} \\
\textbf{Known gaps, all documented} & 4 (TLS, CI/CD, intermittent encryption, CLEAR-as-training) \\
\bottomrule
\end{tabularx}

\vspace{1em}

The project meets \S5.6.1 (70\,\%) and \S5.6.2 (80--85\,\%) in full, and clears five of
seven \S5.6.3 Distinction criteria. The two it does not reach --- Polygon anchoring
and CI/CD --- are placed in Weeks 17--24 by the document's own Tables 5.5 and 5.6.

\textbf{How this was checked.} The PDF was extracted in full and read end to end. All
125 tracked files were opened. The document's API listings (3.1--3.16) were sent
verbatim to the real services and the responses compared field by field. The
benchmarks were re-measured rather than read from \texttt{reports/}. Where a claim and
the code disagreed, the code was treated as the truth and the claim corrected.

% =====================================================================
\section{What changed this session}
% =====================================================================

Ten remediation tasks, then a full re-audit against the PDF that found five more
issues. Fourteen commits on \texttt{fix/phase1-4-remediation}.

\subsection{The blocking defect: small encrypted files produced no response}

\textbf{Where:} \texttt{services/monitor/pipeline.py}

The Monitor flagged them, the ledger recorded them, and nothing acted. The
response gate read only the ML engine's \texttt{threat\_level}, and the behavioural
classifier needs Shannon entropy near 7.995 before it is confident --- which
ciphertext under roughly 40\,KB does not reach through sampling noise. Reproduced
before fixing, five trials per size, in-place encryption keeping the filename:

\begin{center}
\begin{tabular}{@{}llllc@{}}
\toprule
\textbf{Size} & \textbf{Entropy} & \textbf{Monitor verdict} & \textbf{Model p(ransomware)} & \textbf{Response fired} \\
\midrule
4\,KB    & 7.95 & \texttt{suspected\_encryption} & 0.18 & \textbf{No} \\
8\,KB    & 7.98 & \texttt{suspected\_encryption} & 0.04 & \textbf{No} \\
16\,KB   & 7.99 & \texttt{suspected\_encryption} & 0.36 & \textbf{No} \\
32\,KB   & 7.99 & \texttt{suspected\_encryption} & 0.49 & \textbf{No} \\
40\,KB   & 8.00 & \texttt{suspected\_encryption} & 1.00 & Yes \\
64\,KB+  & 8.00 & \texttt{suspected\_encryption} & 1.00 & Yes \\
\bottomrule
\end{tabular}
\end{center}

The ML score now refines the Monitor's verdict instead of overruling it: the
effective threat level is the higher of the two. That also keeps the level
coherent downstream --- the Response service runs network isolation on
\texttt{high}/\texttt{critical} only, so forwarding ``low'' would have triggered a response that
then declined most of its job.

\texttt{pipeline.py} had no direct test coverage, which is how this survived. It now has
21 tests; 7 of them fail against the old gate.

\subsection{The rest}

\begin{xltabular}{\linewidth}{@{}K{0.9cm}Y{0.72}Y{1.28}@{}}
\toprule
\textbf{\#} & \textbf{Finding} & \textbf{What was done} \\
\midrule
\endfirsthead
\toprule
\textbf{\#} & \textbf{Finding} & \textbf{What was done} \\
\midrule
\endhead
\bottomrule
\endfoot

F-2 & Dashboard banner scored a partly fabricated feature vector
    & Four constants came from the document's illustrative example. Fixing them was not
      enough: three more of the model's seven inputs (the byte statistics) were absent
      from the event and being interpolated from entropy --- which is what rendered a
      legitimate ZIP as a threat. The event now carries what was measured, and
      \texttt{handle\_event} derives entropy and statistics from \textbf{one} read instead
      of two. Verified 8/8 across ZIP, text, PDF, 4/8/32/256\,KB encryption and a
      \texttt{.locked} file. \\

F-3 & \texttt{file\_patterns} accepted, echoed, never applied
    & Filtered in \texttt{handle\_event}, so it covers deletes and renames too.
      \texttt{/monitor/status} reports the filter in force. \\

F-4 & A plain \texttt{pytest -q} rewrote three committed benchmark files
    & Writes gated behind \texttt{URDS\_WRITE\_REPORTS=1}. Verified both directions. \\

F-5 & ``Differential Entropy Analysis'' named three times in the spec, not implemented
    & \texttt{EntropyHistory} per path; a rise of $\geq$2.0 bits/byte landing at
      $\geq$7.0 flags replacement. Checked \textbf{before} the container exemption, so it
      catches an encryptor that writes a ZIP header over its ciphertext --- which magic
      bytes alone clear. \\

V-2 & \texttt{/monitor/stop} ignored \texttt{monitor\_id}
    & Accepted and validated; a mismatch is refused with 404. The gateway forwards it
      instead of sending \texttt{\{\}}. \\

V-3 & \texttt{/ledger/log} returned 200
    & Returns 201, per Table 3.2. The proxy passes it through. \\

V-8 & Trained on 19,480 samples, not \S6.1's 50,000
    & Retrained: 35,000 / 7,500 / 7,500, 25,000 per class. \textbf{See \S2 --- the finding
      was not what it first appeared.} \\

D-1 & ``64 features''
    & Actually \textbf{70}. Corrected in 8 places across 4 files; three other ``64''s are
      SHA-256 lengths and were left alone. \\

D-2 & Ledger verification ``$\sim$3.7\,ms''
    & Re-measured properly: \textbf{3.1\,ms median} over 50 warm runs, range 2.1--5.7\,ms.
      The previous audit's 2.3\,ms was a single run near the floor --- the README was
      closer to right than the correction. \\

D-3 & \texttt{attack\_chain\_evidence.txt} headlined ``18/18 checks passed''
    & Actually 16 PASS + 2 SKIP. The generator counted \texttt{None} as a pass; fixed at
      source and in the committed file. \\

D-4 & Table 6.1 describes hardware never used
    & Actual machine recorded against it in \texttt{APPROACH.md} \S10. \\

D-5 & The document gives the hash formula three incompatible ways
    & Footnoted at the design; the code follows Listing 4.7, the document's own reference
      implementation. \\
\end{xltabular}

% =====================================================================
\section{The most important correction: the metrics were stale, not the model}
% =====================================================================

The first pass reported that retraining on 50,000 samples had \textbf{lowered}
accuracy from 0.9774 to 0.9577. That reading was wrong, and the truth matters
more.

The model on disk was \textbf{already} the 50,000-sample model.
\texttt{reports/model\_metrics.json} was a leftover record of an earlier, smaller run
that had since been overwritten. Three independent proofs:

\begin{enumerate}[leftmargin=*,itemsep=0.5em]
  \item \textbf{Training is bit-for-bit deterministic.} Two runs over the same parquet with
        \texttt{random\_state=42} produced \texttt{xgboost\_model.pkl} with the same SHA-256
        (\texttt{78ab2978\ldots d09b3d1}) and the same 1,064,299 bytes. Every metric matched
        to 16 significant digits; only \texttt{trained\_at} differed.
  \item \textbf{SHAP output is unchanged.} \texttt{reports/shap\_feature\_importance.json}
        was committed \emph{before} the retrain, generated from the old model. Re-running
        \texttt{src/shap\_analysis.py} against the new model reproduced it byte-for-byte.
        SHAP values are a deterministic function of tree structure over a seeded sample,
        so identical output means identical trees.
  \item \textbf{The file size never moved.} 1,064,299 bytes before and after. A model fitted
        on 13,636 training rows rather than 35,000 would not land on the same count.
\end{enumerate}

So the deployed classifier has been the 0.9577 model throughout, while
\texttt{model\_metrics.json}, \texttt{/model/metrics}, \texttt{APPROACH.md}, \texttt{FLOW.md},
\texttt{test\_cases.md} and the verification report all cited 0.9774. \textbf{Every document
was overstating the classifier that was actually running, by about two points.}
They now describe the model in the container.

Corroboration: \texttt{reports/ransomware\_specific\_metrics.json} is byte-identical
before and after --- 0.9993 accuracy, 1.0000 precision, 0.9987 recall over 752 real
ransomware-family samples.

\subsection{Current model performance}

\begin{center}
\begin{tabular}{@{}lccccc@{}}
\toprule
\textbf{Model} & \textbf{Accuracy} & \textbf{Precision} & \textbf{Recall} & \textbf{F1} & \textbf{ROC-AUC} \\
\midrule
\textbf{XGBoost (EMBER static PE)} & 0.9577 & 0.9625 & 0.9525 & 0.9575 & 0.9923 \\
Random Forest baseline            & 0.9389 & 0.9490 & 0.9277 & 0.9382 & 0.9859 \\
Behavioural classifier            & 0.8843 & 0.9486 & 0.8127 & 0.8754 & 0.9585 \\
\bottomrule
\end{tabular}
\end{center}

Both EMBER models are trained and evaluated on the same 50,000-row parquet with
the same seed and split, so \textbf{Table 8.1's comparison is apples-to-apples}.
XGBoost beats Random Forest on every metric, which is the relationship Table 8.1
asserts. Every threshold clears: TC-06 wants P/R/F1 above 85\,\%, \S5.6.2 wants ML
above 85\,\%, \S5.6.3 wants above 90\,\% across all metrics.

Ransomware-family detection, measured on the 752 labelled samples in the corpus
(WannaCry 389, GandCrab 343, CryptoLocker 7, Cerber 6, Petya 3, Locky 3,
TeslaCrypt 1): \textbf{99.87\,\% recall, 100\,\% precision}.

% =====================================================================
\section{Chapter 3 --- Architecture \& API, verified line by line}
% =====================================================================

\subsection{Technology stack (Table 3.1)}

\begin{xltabular}{\linewidth}{@{}P{2.2cm}Y{1.00}Y{1.10}Y{0.90}@{}}
\toprule
\textbf{Layer} & \textbf{Document requires} & \textbf{In the repo} & \textbf{Status} \\
\midrule
\endfirsthead
\toprule
\textbf{Layer} & \textbf{Document requires} & \textbf{In the repo} & \textbf{Status} \\
\midrule
\endhead
\bottomrule
\endfoot

Presentation    & Streamlit, Plotly, HTML5/CSS3
                & Streamlit 1.41.1, Plotly 5.24.1, inline CSS & \ok \\
Application     & FastAPI, Python 3.9+, Uvicorn, Pydantic
                & FastAPI 0.115.6, Uvicorn 0.32.1, Pydantic 2.10.4, Python 3.11/3.13 & \ok \\
Business Logic  & XGBoost, watchdog, SHA-256
                & \texttt{xgboost-cpu} 3.4.0, \texttt{watchdog} 6.0.0, \texttt{hashlib.sha256} & \ok \\
Data            & SQLite, file system, Windows VSS
                & \texttt{ledger/database.py}, \texttt{recovery/vss\_manager.py} (WMI) & \ok \\
Security        & JWT, cryptography, \textbf{TLS 1.3}
                & JWT via \texttt{python-jose[cryptography]}; \textbf{no TLS}
                & \warn\ stated in APPROACH \S8 \\
DevOps          & Docker, Compose, \textbf{GitHub Actions}, pytest
                & Docker + Compose + pytest; \textbf{no \texttt{.github/}}
                & \defer\ Table 5.6, Weeks 17--24 \\
\end{xltabular}

\subsection{API contracts --- the document's listings replayed verbatim}

Every listing was sent to the real service and the response compared field by
field. \textbf{53 / 53 checks pass.}

\begin{center}
\begin{tabularx}{\linewidth}{@{}P{3.0cm}K{1.6cm}K{1.4cm}L@{}}
\toprule
\textbf{Service} & \textbf{Checks} & \textbf{Result} & \textbf{Notes} \\
\midrule
Monitor (8001)   & 13 & \ok & Listings 3.1--3.4 exact; responses are supersets \\
ML Engine (8002) & 13 & \ok & Listing 3.5 sent verbatim returns a well-formed 3.6 \\
Ledger (8003)    & 11 & \ok & Listings 3.8/3.9 \textbf{exact match}, no extra fields \\
Response (8004)  & 16 & \ok & Listings 3.10--3.16 all accepted \\
\bottomrule
\end{tabularx}
\end{center}

All 10 documented endpoints exist. Additions (\texttt{/health}, \texttt{/auth/token},
\texttt{/monitor/events}, \texttt{/features}, \texttt{/ledger/entries}, \texttt{/ledger/verify},
\texttt{/ledger/blocks}, \texttt{/analyze}) are required by the orchestration in Listing 4.9,
the dashboard in \S7.1, or TC-05.

\textbf{Hash chain, recomputed from the response alone.}
\texttt{SHA256(timestamp + event\_type + event\_data + previous\_hash)} reproduces
\texttt{current\_hash} exactly; genesis is 64 zeros; block N+1's \texttt{previous\_hash}
equals block N's \texttt{current\_hash}.

\textbf{HTTP status codes (Table 3.2).} 200, 201, 400, 401, 403, 404, 429, 500, 503
all produced. One code is returned that Table 3.2 does not list --- \textbf{409} from
\texttt{/response/terminate} when the guard refuses. It is now declared on the route and
the deviation is recorded in \texttt{APPROACH.md} \S8, with the reasoning: the request is
well-formed and authorised, and 403 already means ``your role is not permitted''
at this gateway.

\textbf{Error envelope (Listing 3.22).} Identical in all six places: gateway, monitor,
ml-engine, ledger, response, recovery.

\textbf{JWT (Listing 3.23).} \texttt{\{sub, role, exp, iat\}} all present, plus \texttt{tier}.

\textbf{Rate limits (Table 3.3).} Free 60/100, Premium 300/500, Enterprise 1000/2000 ---
exact transcription.

\subsection{Data models (\S3.5)}

\begin{center}
\begin{tabularx}{\linewidth}{@{}P{3.6cm}L@{}}
\toprule
\textbf{Model} & \textbf{Status} \\
\midrule
\texttt{FileEvent} (3.17)   & All fields present; \texttt{hash\_md5} $\rightarrow$ \texttt{file\_hash} holding SHA-256 (justified in APPROACH \S8) \\
\texttt{Prediction} (3.18)  & Exact \\
\texttt{LedgerBlock} (3.19) & Exact, including \texttt{blockchain\_anchor} present and always null \\
\bottomrule
\end{tabularx}
\end{center}

\subsection{Sequence diagrams (\S3.6)}

\begin{itemize}[leftmargin=*,itemsep=0.4em]
  \item \textbf{Figure 3.4 (attack):} implemented exactly --- modification $\rightarrow$
        entropy $\rightarrow$ features $\rightarrow$ RANSOMWARE $\rightarrow$ trigger
        $\rightarrow$ kill $\rightarrow$ alert $\rightarrow$ isolate.
  \item \textbf{Figure 3.3 (benign):} diverges deliberately. Benign events are recorded but
        do not fan out to ML and the ledger, which is what keeps CPU near 1\,\%. The
        documented flow is still reachable via \texttt{POST /analyze}. Stated in
        APPROACH \S8.
\end{itemize}

\subsection{Deployment (\S3.7.1, Listing 3.20)}

\texttt{docker-compose.yml} is a strict superset. \texttt{docker compose config} validates.
Every item in Listing 3.20 is present including the exact \texttt{cpus: '2' / memory: 2G}
limits; ports are bound to \texttt{127.0.0.1} for the four backends (hardening), and
\texttt{depends\_on} is upgraded to \texttt{condition: service\_healthy}.

% =====================================================================
\section{Chapters 4--5 --- deliverables, test cases, criteria}
% =====================================================================

\subsection{Weekly deliverables (Tables 4.1--4.4)}

All Weeks 1--16 rows met. Notable: entropy calculator tested against known values
(uniform = 8.0, single byte = 0.0, two symbols = 1.0); PE pipeline extracts \textbf{70}
features against a ``50+'' requirement; VSS snapshots every 6 hours
(\texttt{DEFAULT\_INTERVAL\_HOURS = 6}); backup creation measured at 2.8\,s against a 30\,s
target. Several Weeks 17--24 rows are already met early (false positives $<$5\,\%,
model $>$90\,\%, SHAP documented).

\subsection{Test cases (Table 5.8)}

\begin{center}
\begin{tabularx}{\linewidth}{@{}P{1.5cm}Y{0.85}Y{1.15}@{}}
\toprule
\textbf{ID} & \textbf{Scenario} & \textbf{Verified} \\
\midrule
TC-01 & Known ransomware, $<$5 files encrypted   & \ok\ detected after \textbf{1 file} \\
TC-02 & Zero-day / behavioural                   & \ok \\
TC-03 & Legitimate ZIP, no alert                 & \ok\ 0/40 false positives \\
TC-04 & Recovery from backup                     & \ok\ native \\
TC-05 & Audit-log tampering detected             & \ok\ exact block identified \\
TC-06 & P/R/F1 all $>$85\,\%                      & \ok\ 0.9625 / 0.9525 / 0.9575 \\
TC-07 & Terminated within 2\,s                   & \ok\ 125.6\,ms \\
TC-08 & CPU $<$15\,\%, RAM $<$500\,MB             & \ok\ 0.96\,\% / 70.3\,MB --- \textbf{now \texttt{tc08}-named} \\
TC-09 & Dashboard alert within 1\,s              & \ok\ \textbf{10.5\,ms --- this had no automated test before this session} \\
TC-10 & 401 + audit entry                        & \ok\ all three parts \\
TC-11 & Simultaneous attacks                     & \ok\ 12 concurrent detections, 8 concurrent kills \\
TC-12 & Polygon anchoring \textit{(if implemented)} & \defer\ Weeks 17--24 \\
\bottomrule
\end{tabularx}
\end{center}

\subsection{Testing pyramid (\S5.5.1)}

\begin{center}
\begin{tabular}{@{}lcc@{}}
\toprule
\textbf{Tier} & \textbf{Document target} & \textbf{Actual} \\
\midrule
Unit        & 100+ & \textbf{165} \\
Integration & 20   & \textbf{177} \\
E2E         & 5    & \textbf{25}  \\
\bottomrule
\end{tabular}
\end{center}

\subsection{Performance benchmarks (Table 5.9) --- all re-measured}

\begin{center}
\begin{tabularx}{\linewidth}{@{}P{4.0cm}P{2.0cm}L@{}}
\toprule
\textbf{Metric} & \textbf{Target} & \textbf{Measured} \\
\midrule
Detection latency    & $<$100\,ms  & \textbf{23.7\,ms} p95 (40 files, 4\,KB--2\,MB) \\
Response time        & $<$2\,s     & \textbf{125.6\,ms} \\
False positive rate  & $<$5\,\%     & \textbf{0\,\%} (0/40, 32 deliberately high-entropy) \\
ML inference         & $<$100\,ms  & \textbf{1.95\,ms} p95 end-to-end, 0.44\,ms model-only \\
API response p95     & $<$200\,ms  & \textbf{3.22\,ms} over 1000 requests \\
System CPU           & $<$15\,\%    & \textbf{0.96\,\%} of 14 cores \\
System RAM           & $<$500\,MB  & \textbf{70.3\,MB} peak \\
File recovery        & 100\,\%      & \textbf{100\,\%} native \\
Ledger verification  & $<$50\,ms   & \textbf{3.1\,ms} median / 1000 blocks \\
Dashboard latency    & $<$1\,s     & \textbf{10.5\,ms} to queryable; 439\,ms end-to-end in Compose \\
\bottomrule
\end{tabularx}
\end{center}

\subsection{Success criteria (\S5.6)}

\begin{center}
\begin{tabularx}{\linewidth}{@{}P{2.6cm}Y{1.20}Y{0.80}@{}}
\toprule
\textbf{Tier} & \textbf{Criterion} & \textbf{Met} \\
\midrule
\S5.6.1 (70\,\%)    & Monitoring + entropy                    & \ok \\
                   & ML $>$70\,\%                             & \ok\ 95.8\,\% \\
                   & Stops $\geq$1 simulator                 & \ok \\
                   & Hash-chain audit trail                  & \ok \\
                   & Basic recovery                          & \ok \\
\addlinespace
\S5.6.2 (80--85\,\%)& ML $>$85\,\% (P/R/F1)                    & \ok \\
                   & \textbf{3+ ransomware simulators}       & \ok\ \textbf{now three} --- see \S5 \\
                   & False positives $<$5\,\%                 & \ok\ 0\,\% \\
                   & Response $<$2\,s                        & \ok\ 125.6\,ms \\
                   & 8--10 test cases                        & \ok\ 11 \\
                   & Dashboard with real-time alerts         & \ok \\
\addlinespace
\S5.6.3 (90\,\%+)   & ML $>$90\,\% all metrics                 & \ok \\
                   & Blockchain anchoring                    & \defer\ Weeks 17--24 \\
                   & Professional dashboard                  & \ok \\
                   & 12--15 test cases                       & \warn\ 11 of 12 (TC-12 gates the 12th) \\
                   & CI/CD pipeline                          & \defer\ Weeks 17--24 \\
                   & SHAP documented                         & \ok \\
                   & Research paper draft                    & \defer\ Weeks 25--32 \\
\bottomrule
\end{tabularx}
\end{center}

\subsection{Risk register (Table 5.7)}

\begin{center}
\begin{tabularx}{\linewidth}{@{}P{3.2cm}Y{0.92}Y{1.08}@{}}
\toprule
\textbf{Risk} & \textbf{Document's mitigation} & \textbf{In the repo} \\
\midrule
ML underperforms      & Random Forest fallback                        & \ok\ trained, fills Table 8.1 \\
Integration issues    & API contracts by Week 4, Compose              & \ok\ bidirectional parity tests \\
High false positives  & Differential entropy, magic bytes, whitelist  & \ok\ \ok\ / \no\ whitelist not built (FP rate is 0\,\%) \\
Blockchain costs      & Testnet, hash chain fallback                  & \ok\ fallback is the shipped path \\
Dataset access        & Smaller subsets                               & \ok\ \texttt{--per-class} supported \\
\bottomrule
\end{tabularx}
\end{center}

% =====================================================================
\section{Ransomware simulators --- \S5.6.2's ``3+'' requirement}
% =====================================================================

Previously one simulator run in four configurations, which is one behaviour. It
now imitates four families that differ in \textbf{what the watcher sees}, not in the
payload. Measured, not assumed:

\begin{center}
\begin{tabularx}{\linewidth}{@{}P{2.2cm}LK{2.2cm}@{}}
\toprule
\textbf{Family} & \textbf{Behaviour} & \textbf{Detected} \\
\midrule
\texttt{locker}  & Rewrite in place, append \texttt{.locked}                                          & \textbf{8/8} \\
\texttt{silent}  & Rewrite in place, keep the name --- no extension signal                            & \textbf{8/8} \\
\texttt{copycat} & Write a new encrypted file, delete the original (\texttt{created}+\texttt{deleted}) & \textbf{8/8} \\
\texttt{partial} & Scramble the leading quarter (LockBit 3 / BlackCat)                                & \textbf{0/8} \\
\bottomrule
\end{tabularx}
\end{center}

Three families detected satisfies \S5.6.2. All four round-trip through
\texttt{--restore}, which is what keeps them safe to run.

\textbf{\texttt{partial} is a genuine blind spot and is asserted as one.} Intermittent
encryption leaves whole-file entropy near \textbf{5.2 bits/byte} against a 7.5
threshold, so it reads as an ordinary edit. Catching it needs per-block entropy
rather than a whole-file average, which would also flag the compressed blocks
inside any \texttt{.docx} or PDF --- and the 0\,\% false-positive rate is a graded criterion
while intermittent encryption is not mentioned anywhere in the reference
document. The gap is captured by
\texttt{test\_\allowbreak partial\_\allowbreak encryption\_\allowbreak is\_\allowbreak a\_\allowbreak known\_\allowbreak blind\_\allowbreak spot},
so it is visible in the suite rather than absent from it, and that test fails loudly
if detection ever improves.

% =====================================================================
\section{Chapter 6 --- experimental methodology}
% =====================================================================

\begin{center}
\begin{tabularx}{\linewidth}{@{}P{3.2cm}Y{1.05}Y{0.95}P{2.9cm}@{}}
\toprule
\textbf{\S6.1 requirement} & \textbf{Document} & \textbf{Actual} & \textbf{Status} \\
\midrule
Total samples    & 50,000 PE files                        & \textbf{50,000}            & \ok \\
Class balance    & 25,000 / 25,000                        & \textbf{25,000 / 25,000}   & \ok \\
Split ratio      & 70 / 15 / 15                           & \textbf{35,000 / 7,500 / 7,500} & \ok\ exact \\
Training source  & EMBER \textbf{+ CLEAR behavioral logs} & EMBER only; CLEAR used for EDA & \warn\ stated in APPROACH \S8 \\
\bottomrule
\end{tabularx}
\end{center}

\S6.2's five metrics are all computed and reported. \S6.3's Table 6.1 describes an
i7-12700K / 32\,GB / Python 3.9.13 that was never used; the actual machine
(Core Ultra 5 225H, 14 cores, 15\,GB, Windows 11 build 26200, Python 3.13.9) is
recorded against it in \texttt{APPROACH.md} \S10. Every target is met with room on this
hardware, so no conclusion depends on the difference.

% =====================================================================
\section{Known gaps, all documented}
% =====================================================================

Each is stated in \texttt{APPROACH.md} \S8 or \S9 rather than left to be discovered.

\begin{xltabular}{\linewidth}{@{}Y{0.62}Y{1.38}@{}}
\toprule
\textbf{Gap} & \textbf{Why it stands} \\
\midrule
\endfirsthead
\toprule
\textbf{Gap} & \textbf{Why it stands} \\
\midrule
\endhead
\bottomrule
\endfoot

\textbf{No TLS 1.3 between services} (Table 3.1)
  & Real gap. Mitigated by binding the four backends to loopback so unauthenticated
    traffic never crosses a network boundary. \S3.7.3 scopes TLS to production. \\

\textbf{Intermittent encryption not detected}
  & Needs per-block entropy; the false-positive cost is not worth it for Weeks 1--16,
    and the spec never asks for it. \\

\textbf{CLEAR is EDA-only, not training input} (\S6.1)
  & Folding it in means retraining and restating every cited figure. Deferred
    deliberately, not half-done. \\

\textbf{SMOTE omitted} (p.\,40)
  & The corpus is balanced 25,000/25,000, so \texttt{scale\_pos\_weight} computes to
    $\sim$1.0 and SMOTE has no minority to oversample. Implementing it would be a
    no-op that looked like a safeguard. \\

\textbf{Whitelist / training mode} (Table 5.7)
  & Two of four listed FP mitigations are built; the measured FP rate is 0\,\%, so the
    remaining two are unnecessary for now. \S8.5 lists whitelisting as future work. \\

\textbf{Behavioral fingerprinting is static} (\S1.4)
  & API-call features come from the PE import table grouped by behaviour class, not
    from runtime call sequences. Runtime attribution needs ETW/eBPF --- Phase 5. \\

\textbf{Process attribution is null}
  & Watchdog reports \emph{what} changed, never \emph{who}. Reporting the monitor's own
    PID would name an unrelated process to the response container. \\

\textbf{VSS deletion protection}
  & \S2.5 reviews it but no Weeks 1--16 deliverable lists it. Correctly out of scope;
    named in the limitations. \\

\textbf{CI/CD, Polygon/TC-12, 95\,\% coverage, SHAP per-prediction, research paper}
  & All Weeks 17--32 by Tables 5.4--5.6. \\
\end{xltabular}

\vspace{0.8em}

Minor, non-blocking: the directory is \texttt{services/ml-engine} where Appendix A.1
writes \texttt{ml\_engine} and Listing 3.20 writes \texttt{ml} --- the document is internally
inconsistent, and the Compose \emph{service} name is \texttt{ml\_engine}, so the DNS name in
Listing 3.20's \texttt{ML\_URL=http://ml\_engine:8002} resolves exactly as documented.
Entry points are \texttt{app.py}/\texttt{main.py} rather than A.1's \texttt{main.py}/\texttt{api.py}.

% =====================================================================
\section{Reproducing this}
% =====================================================================

\begin{lstlisting}
for svc in gateway ledger monitor ml-engine response; do (cd services/$svc && python -m pytest -q); done
\end{lstlisting}

Expect \textbf{371 passed, 2 skipped, 0 failed}. The two skips are legitimate:
\texttt{psutil.terminate()} maps to \texttt{TerminateProcess} on Windows, which no process can
ignore, so the SIGTERM-escalation tests assert a POSIX guarantee with no Windows
equivalent; a Windows-specific test covers the same ground.

Benchmarks, with measured values printed:

\begin{lstlisting}
cd services/monitor && python -m pytest -m benchmark -q -s
\end{lstlisting}

A plain test run leaves \texttt{reports/} untouched. To refresh the committed evidence
deliberately:

\begin{lstlisting}
URDS_WRITE_REPORTS=1 python -m pytest -m benchmark -q -s
\end{lstlisting}

% =====================================================================
\section{Commits}
% =====================================================================

\begin{center}
\begin{tabularx}{\linewidth}{@{}P{2.3cm}L@{}}
\toprule
\textbf{Commit} & \textbf{Task} \\
\midrule
\texttt{c4b9ed1} & \texttt{fix(monitor)}: let the Monitor's verdict set a floor on threat level (F-1) \\
\texttt{3707b22} & \texttt{fix(dashboard)}: score the file that was actually detected (F-2) \\
\texttt{b3642d7} & \texttt{fix(monitor)}: apply \texttt{file\_patterns} instead of only echoing it (F-3) \\
\texttt{4a3fdbb} & \texttt{test}: record benchmark numbers only when asked (F-4) \\
\texttt{173c5c0} & \texttt{feat(monitor)}: implement differential entropy analysis (F-5) \\
\texttt{9db64a1} & \texttt{fix}: three spec-exactness corrections (V-2, V-3, D-1) \\
\texttt{7368cc3} & \texttt{feat(ml)}: retrain on the 50,000 samples the specification calls for (V-8) \\
\texttt{9b9c11c} & \texttt{docs}: correct the numbers, state the deviations, mark the dead code \\
\texttt{d3aeda3} & \texttt{fix}: close the gaps a full re-audit against the reference PDF found \\
\bottomrule
\end{tabularx}
\end{center}

\end{document}
